\begin{theorem}[High stationary classes: unramified and stable cases]
\label{mod:parameters-highunramifiedstable}
Let $K/F$ be a finite valuative extension of nonarchimedean local fields and
let
\[
 \chi_K=\chi_F\circ N_{K/F},\qquad
 \psi_K=\psi_F\circ\operatorname{Tr}_{K/F}.
\]
Write the separately stored exact multiplicative conductors as
\[
 m_F=2d_F+\varepsilon_F,\qquad
 m_K=2d_K+\varepsilon_K,\qquad
 \varepsilon_F,\varepsilon_K\leq1,\qquad m_F,m_K>1,
\]
and put $q_F=d_F+\varepsilon_F$ and
$q_K=d_K+\varepsilon_K$.  Given a downstairs denominator $\gamma_F$ with
\[
 v_F(\gamma_F)=m_F+n_F(\psi_F),
\]
the theorems below use the ordered pair
\[
 (\gamma_F,\gamma_K)
   =(\gamma_F,\operatorname{algebraMap}_{F,K}(\gamma_F)).
\]
In every case Lean first proves that $\gamma_K$ has the required upstairs
valuation $m_K+n_K(\psi_K)$.

Let
\[
 \beta_F\in\mathcal O_F/\mathfrak p_F^{d_F},\qquad
 \beta_K\in\mathcal O_K/\mathfrak p_K^{d_K}
\]
be the stationary numerator classes constructed at the exact depths $q_F$
and $q_K$.  Whenever the truncated norm is trace at these depths, combining
the stationary-class transport of
Corollary~\ref{mod:parameters-stationaryclassundernorm} with the
denominator-sensitive adjoint gives
\[
 \boxed{\quad
 \beta_K=
 P^*_{K/F;\gamma_F,\gamma_K}(\beta_F)
 =\iota^{d_F,d_K}_{\gamma_F,\gamma_K}(\beta_F)
 \quad}
 \qquad\text{in }\mathcal O_K/\mathfrak p_K^{d_K}.
\]
Here $P$ points from the upstairs variable quotient to the downstairs
variable quotient, $P^*$ points from the downstairs numerator quotient to
the upstairs numerator quotient, and
$\iota^{d_F,d_K}_{\gamma_F,\gamma_K}$ is Lean's
\texttt{denominatorScaledAlgebraMap}.  Thus the right side is a quotient
class obtained from the algebra map with this ordered denominator pair; it
is not a chosen element of $K$.

\paragraph{Unramified row.}
Assume $K/F$ is prime cyclic and $e(K/F)=1$, so
$f(K/F)=[K:F]$.  The factor $e=1$ controls the common denominator and the
factor $f=[K:F]$ is retained in the terminal norm bound.  Lean proves the exact
preservation identities
\[
 m_K=m_F,\qquad n_K(\psi_K)=n_F(\psi_F),
\]
and uses the same decomposition $(d_K,\varepsilon_K)=(d_F,\varepsilon_F)$.
It constructs the three norm-filtration inclusions at the exact pairs
\[
 (q_K,q_F)=(q_F,q_F),\qquad (m_K,m_F)=(m_F,m_F),
 \qquad(d_K,d_F)=(d_F,d_F).
\]
For $x\in\mathfrak p_K^{q_F}$ the degree-$j$ elementary symmetric term has
downstairs valuation at least $j q_F$; hence every term with $j\geq2$
vanishes modulo $\mathfrak p_F^{m_F}$ because $m_F\leq2q_F$.  Therefore
\[
 P:[x]\longmapsto[N_{K/F}(1+x)-1]
   =\operatorname{Tr}_{K/F}
\]
on
$\mathfrak p_K^{q_F}/\mathfrak p_K^{m_F}
 \to\mathfrak p_F^{q_F}/\mathfrak p_F^{m_F}$, and the displayed quotient
identity is precisely the specialization $P^*(\beta_F)=\beta_F$ in
$\mathcal O_K/\mathfrak p_K^{d_F}$.

\paragraph{Stable odd-degree ramified row.}
Assume $K/F$ is prime cyclic, $f(K/F)=1$, and supply the lower break $t$ and
the generating integral uniformizer required by the ramification lemmas.
Put
\[
 \ell=[K:F]=e(K/F),\qquad T=t+1,\qquad
 D=(\ell-1)T,
\]
assume $\ell$ is odd and retain the manuscript's exact stable inequality
$T\leq d_F$.  No further tame/wild disjunction is assumed.  Lean splits
internally on $t=0$; in the positive-break branch it derives
$\operatorname{char}(k_F)=\ell$ from the supplied lower-break and generating
uniformizer data.
The conductor results of Theorem~\ref{mod:ramification-pullbackconductors}
give
\[
 m_K+D=\ell m_F,
 \qquad
 m_K+n_K(\psi_K)=\ell\bigl(m_F+n_F(\psi_F)\bigr),
\]
so the common denominator is admissible.  Writing $\ell=2k+1$, the Lean
arithmetic proof derives, without truncated-subtraction ambiguity,
\[
 \varepsilon_K=\varepsilon_F,\qquad
 d_K=d_F+k(m_F-T),\qquad
 q_K=q_F+k(m_F-T).
\]
In particular it proves
\[
 q_K\geq m_F,\qquad d_K\leq\ell d_F,\qquad
 \ell q_F\leq q_K+D,\qquad
 \ell m_F=m_K+D,
\]
as well as the coefficient-depth bound needed for
$N(U_K^{d_K})\subseteq U_F^{d_F}$.  Lean constructs, rather than assumes,
all three exact-precision containments
\[
 N(U_K^{q_K})\subseteq U_F^{q_F},\qquad
 N(U_K^{m_K})\subseteq U_F^{m_F},\qquad
 N(U_K^{d_K})\subseteq U_F^{d_F},
\]
and the trace containments at $(q_K,q_F)$ and $(m_K,m_F)$.

For $t=0$ it applies the totally ramified tame elementary-symmetric bound;
for $t>0$ it applies the wild bound with
$\operatorname{char}(k_F)=\ell$.  The shared numerical estimate
\[
 \ell m_F\leq2q_K+D
\]
kills every intermediate term modulo $\mathfrak p_F^{m_F}$.  The terminal
norm term has $v_F(Nx)=f(K/F)v_K(x)=v_K(x)\geq q_K\geq m_F$.
Consequently $P=\operatorname{Tr}_{K/F}$ on the exact variable quotients,
and Lean returns a proof of $d_K\leq e(K/F)d_F$ together with the displayed
stationary numerator quotient equality.

\paragraph{Tamely ramified quadratic row.}
Assume further
\[
 \ell=e(K/F)=2,\qquad f(K/F)=1,\qquad t=0,\qquad
 \operatorname{char}(k_F)\ne2.
\]
From the actual pullback conductor, rather than an independently supplied
upstairs depth, Lean proves
\[
 m_K+1=2m_F,\qquad
 m_K=2(m_F-1)+1,\qquad
 d_K=m_F-1,\qquad \varepsilon_K=1,\qquad q_K=m_F.
\]
Thus the truncated norm has the exact source and target
\[
 P:\mathfrak p_K^{m_F}/\mathfrak p_K^{2m_F-1}
   \longrightarrow
   \mathfrak p_F^{d_F+\varepsilon_F}/\mathfrak p_F^{m_F}.
\]
The intermediate range is empty in degree two, while
$v_F(Nx)=v_K(x)\geq m_F$ kills the quadratic terminal term.  Hence the
quadratic identity $N(1+x)-1=\operatorname{Tr}(x)+N(x)$ specializes to
$P=\operatorname{Tr}_{K/F}$ on these exact quotients.  Lean also proves
$m_F-1\leq2d_F$ and the common-denominator identity
\[
 v_K(\gamma_K)=2\bigl(m_F+n_F(\psi_F)\bigr)
               =m_K+n_K(\psi_K),
\]
then concludes the quotient-class equality in
$\mathcal O_K/\mathfrak p_K^{m_F-1}$.  No equality of representatives in
$K$, and no choice of a quadratic uniformizer, is asserted.

\LeanFile{LanglandsFirstMainLemma/Parameters/HighUnramifiedStable.lean}
\lean{LanglandsFirstMainLemma.highParameter_unramified_conductor_eq}
\lean{LanglandsFirstMainLemma.highParameter_unramified_additiveConductor_eq}
\lean{LanglandsFirstMainLemma.highParameter_unramified_sourceDecomposition}
\lean{LanglandsFirstMainLemma.highParameter_unramified_commonDenominator}
\lean{LanglandsFirstMainLemma.highParameter_unramified}
\lean{LanglandsFirstMainLemma.highParameter_ramified_conductor_eq}
\lean{LanglandsFirstMainLemma.highParameter_ramified_conductor_relation}
\lean{LanglandsFirstMainLemma.highParameter_ramified_commonDenominator}
\lean{LanglandsFirstMainLemma.highParameter_stableOdd}
\lean{LanglandsFirstMainLemma.highParameter_tameQuadratic_conductor_relation}
\lean{LanglandsFirstMainLemma.highParameter_tameQuadratic_sourceDecomposition}
\lean{LanglandsFirstMainLemma.highParameter_tameQuadratic_variableDepth}
\lean{LanglandsFirstMainLemma.highParameter_tameQuadratic}
\leanok
\SourceText{Proposition prop:high-parameter-table, parts (a), (b), and the
tame quadratic part.}
\uses{mod:parameters-stationaryclassundernorm,
mod:ramification-pullbackconductors,
mod:ramification-unramifiedcompatibility}
\FormalizationContract
The three principal results are equalities in literal coefficient lattice
quotients.  Their right sides are the denominator-sensitive algebra-map
images for the common ordered denominator pair.  The unramified theorem uses
depths $(q_F,m_F,d_F)$ on both sides; the stable odd theorem uses the actual
upstairs depths $(q_K,m_K,d_K)$ determined by norm pullback and assumes
exactly $t+1\leq d_F$; the tame quadratic theorem derives
$(q_K,m_K,d_K)=(m_F,2m_F-1,m_F-1)$.  Norm maps from upstairs filtrations to
downstairs filtrations, while its adjoint maps downstairs numerator classes
to upstairs numerator classes.  The proofs use the exact trace
specialization only after constructing the required conductor, variable,
and ambiguity precision and after proving the common denominator admissible.
They do not select stationary representatives and do not import any later
high-parameter, phase-reduction, or case theorem.
\end{theorem}
